\section{Introduction}
\label{sec:intro}

Assistants built on large language models (LLMs) increasingly act on a user's behalf: they open pages, fill forms and read results. Several of the web and GUI agents reviewed in Section~\ref{sec:related} send each observation, often a full screenshot, to a hosted multimodal model~\cite{he2024webvoyager,zheng2024seeact,zhang2025ufo,zhang2025appagent}. For a personal assistant this is a poor fit. The screen of a logged-in user carries names, messages and account data, and leakage of personal information from language models has been measured even after scrubbing~\cite{lukas2023pii}. Few of these systems keep memory across sessions, and most operate either in a browser or in one other environment rather than across the applications a person uses (Table~\ref{tab:related}).

A second problem is independent of where the model runs. An agent can believe it has finished when it has not. We call this \emph{premature completion}: the agent emits a terminal ``done'' while the environment does not satisfy the user's goal. A typical case is a search request in which the agent opens the search engine, sees the page load without error, and stops before any query is issued. Every individual browser call succeeded; the task did not. Benchmarks that check final environment state show how wide this gap is~\cite{zhou2024webarena,xie2024osworld}, and self-assessment by the model is not a reliable remedy, since intrinsic self-correction without external feedback does not improve, and can degrade, reasoning accuracy~\cite{huang2024cannot}.

This paper describes Agentic Pilot, an open-source agent that runs its language and vision models locally through Ollama and drives a Chromium browser through Playwright. Its central design rule is that a task may be reported complete only when a check on the observed environment state passes, and that the inputs to this check are written to disk as evidence. The agent is structured as an explicit state machine in LangGraph, so that observation, planning, execution, verification and recovery are separate nodes whose transitions can be inspected and logged.

We make the following contributions:
\begin{enumerate}
\item An evidence-gated execution model that separates action dispatch from task completion, formalized as a completion predicate over observations, together with a \emph{false completion rate} (FCR) metric for measuring premature completion (Section~\ref{sec:formal}).
\item The architecture and implementation of a local agent that combines this gate with role-based routing across local models, DOM-first grounding with a vision-model fallback, CAPTCHA detection that halts rather than bypasses, a task-level human approval gate, and persistent per-step evidence (Sections~\ref{sec:arch}--\ref{sec:impl}).
\item A controlled, reproducible study of the completion gate on 35 constructed browser end states, which quantifies both the premature completions it prevents and the ones it misses (Section~\ref{sec:results}).
\item A design for extending the framework from the browser to the local system (files, processes and native applications) with accessibility-tree-first grounding, typed desktop verification predicates and deterministic permission tiers; we state precisely which parts of it exist in code today (Section~\ref{sec:desktop}).
\item A code-level audit of which mechanisms are active on the live execution path, and an end-to-end evaluation protocol with ablations for the author's hardware (Sections~\ref{sec:impl} and~\ref{sec:setup}).
\end{enumerate}

Section~\ref{sec:related} reviews related work. Section~\ref{sec:arch} presents the architecture and Section~\ref{sec:method} the methodology, including the formal model and the local-system extension. Section~\ref{sec:impl} covers the implementation and its status, Section~\ref{sec:setup} the experimental setup and Section~\ref{sec:results} the results. Sections~\ref{sec:security} and~\ref{sec:limits} discuss security and limitations, and Section~\ref{sec:conclusion} concludes.
